\subsection*{Assignment 6 -- Data Uploading with Python}

L’Assignment~6 richiedeva lo sviluppo di un programma Python in grado di
popolare il database del Data Warehouse (\textit{Group\_ID\_8\_DB}),
rispettando lo schema e le relazioni definite nell’Assignment~4 e utilizzando
i dati preparati negli assignment precedenti.

Poiché il caricamento completo dei dati può risultare oneroso in termini di
tempo e risorse, è stato richiesto di progettare il codice in modo appropriato,
tenendo conto sia delle prestazioni sia della stabilità del processo.

Il caricamento dei dati è stato implementato tramite uno script Python che
utilizza la libreria \texttt{pyodbc} per stabilire una connessione a
\textit{SQL Server}. La connessione viene effettuata specificando server,
database e credenziali, consentendo l’esecuzione diretta di istruzioni SQL
dal codice Python. Una volta stabilita la connessione, viene creato un cursore
che permette di eseguire operazioni di inserimento sulle tabelle del
Data Warehouse.

Per garantire il rispetto dei vincoli di integrità referenziale, le tabelle
vengono caricate seguendo un ordine prestabilito, coerente con le relazioni
di \textit{foreign key} definite nello schema:
\begin{itemize}
    \item dimensioni indipendenti;
    \item dimensioni che dipendono da altre dimensioni;
    \item fact table.
\end{itemize}

In particolare, le dimensioni vengono caricate prima della fact table
\textit{FactParticipation}, assicurando che tutte le chiavi esterne siano già
presenti nel database al momento dell’inserimento dei fatti.

Prima di avviare il caricamento, lo script verifica se una tabella è già
popolata tramite una semplice query di conteggio. Se una tabella risulta non
vuota, il caricamento viene saltato, evitando duplicazioni accidentali dei
dati. Questo controllo rende lo script riutilizzabile e consente di rilanciare
il processo senza dover ripulire manualmente il database.

I dati vengono letti dai file CSV prodotti nell’Assignment~5. Per ogni tabella,
lo script:
\begin{itemize}
    \item legge l’intestazione del file per determinare le colonne;
    \item costruisce dinamicamente la query di inserimento;
    \item effettua una conversione esplicita dei tipi di dato per alcuni
    attributi critici, in particolare:
    \begin{itemize}
        \item \textit{Streams1Month}, convertito in valore numerico;
        \item \textit{IsPrimary}, convertito in valore booleano/integer;
    \end{itemize}
    \item gestisce correttamente i valori mancanti, convertendoli in
    \texttt{NULL}.
\end{itemize}

Poiché il caricamento completo dei dati può richiedere tempi significativi,
l’inserimento delle tuple è stato progettato in modalità \textit{batch}. I
record vengono accumulati in gruppi di dimensione fissa e inseriti nel database
tramite operazioni di inserimento multiplo, seguite da un commit esplicito.

Questa scelta consente di ridurre il numero di comunicazioni tra Python e il
database, migliorare le prestazioni complessive del caricamento e rendere il
processo più stabile in presenza di grandi volumi di dati.

Lo script presentato rappresenta la versione finale del processo di
caricamento: durante lo sviluppo sono stati effettuati diversi test e lanci
intermedi per verificare la connessione al database e il corretto inserimento
dei dati, prima di consolidare la soluzione a batch adottata.

L’Assignment~6 conclude quindi il processo di costruzione del Data Warehouse,
collegando la fase di preparazione dei dati alla loro effettiva persistenza nel
database in modo automatico, coerente e ripetibile.
